% Section 3, part 3 of 8: the commutant, lifting, the scalar equation.

\begin{lem}[The commutant]
\label{lem:gau-comm}
Let $F$ be a field of characteristic $0$, $\tilde b_0,\dots,\tilde b_4\in F^\times$, $\tilde A=\sum_a\tilde b_aE_{a,a+1}$, and define $\mathfrak{gl}^{(m)}$, $\Bprod^{(m)}_a$, $\Bprod$ and $\mathcal Y$ over $F$ as above.
\begin{enumerate}
\item[(a)] $\tilde A^m=\sum_a\Bprod^{(m)}_aE_{a,a+m}$ for $0\leq m\leq5$, so $\tilde A^5=\Bprod\,1$; $\mathcal Y=\sum_a\tilde b_a^{-1}E_{a+1,a}$, $\mathcal Y^j=\sum_a(\Bprod^{(j)}_a)^{-1}E_{a+j,a}$ for $0\leq j\leq5$, and $\tilde A^m=\Bprod\,\mathcal Y^{5-m}$ for $1\leq m\leq5$; $\det(x1-\tilde A)=x^5-\Bprod$.
\item[(b)] For $X=\sum_a\xi_aE_{a,a+m}\in\mathfrak{gl}^{(m)}$, $[\tilde A,X]=\sum_a(\tilde b_a\xi_{a+1}-\xi_a\tilde b_{a+m})E_{a,a+m+1}\in\mathfrak{gl}^{(m+1)}$.
\item[(c)] For every $m$, $\ker(\operatorname{ad}\tilde A)\cap\mathfrak{gl}^{(m)}=F\tilde A^m$; hence $\ker(\operatorname{ad}\tilde A)=F[\tilde A]=F[\mathcal Y]=\bigoplus_{j=0}^4F\mathcal Y^j$, of dimension $5$.
\item[(d)] For $m\in\{1,\dots,5\}$ and $X=\sum_a\xi_aE_{a,a+m}\in\mathfrak{gl}^{(m)}$ put $\lambda_m(X):=\sum_a\xi_a/\Bprod^{(m)}_a$ (a local letter). Then $\lambda_m(\tilde A^m)=5$, $\operatorname{ad}\tilde A(\mathfrak{gl}^{(m-1)})=\ker\lambda_m$, $\mathfrak{gl}^{(m)}=F\tilde A^m\oplus\ker\lambda_m$, and $\mathfrak{gl}_5(F)=F[\mathcal Y]\oplus\operatorname{im}(\operatorname{ad}\tilde A)$.
\item[(e)] The projection $\operatorname{pr}$ onto $F[\mathcal Y]$ along $\operatorname{im}(\operatorname{ad}\tilde A)$ is $F$-linear and preserves the shift; for $X\in\mathfrak{gl}^{(m)}$, $1\leq m\leq5$,
\[
\operatorname{pr}(X)=\tfrac15\lambda_m(X)\,\tilde A^m=\tfrac15\Bigl[\sum_a\xi_a\Bprod^{(5-m)}_{a+m}\Bigr]\mathcal Y^{5-m}.
\]
Equivalently, $\operatorname{pr}(X)$ is the unique element of $F[\mathcal Y]$ with $\tr(\operatorname{pr}(X)Z)=\tr(XZ)$ for all $Z\in F[\mathcal Y]$.
\item[(f)] Let $F=\Bbbk\lser u$ and let $\tilde A$ have the one-pole shape. Then $\mathcal Y\in\mathfrak{gl}_5(\Bbbk\fps u)$; an element $\sum_{j=0}^4w_j\mathcal Y^j$ ($w_j\in\Bbbk\lser u$) lies in $\mathfrak{gl}_5(\Bbbk\fps u)$ if and only if every $w_j$ lies in $\Bbbk\fps u$; the commutant of $\tilde A$ in $\mathfrak{gl}_5(\Bbbk\fps u)$ is the ring $\Bbbk\fps u[\mathcal Y]:=\bigoplus_{j=0}^4\Bbbk\fps u\,\mathcal Y^j$; and for $X\in\mathfrak{gl}_5(\Bbbk\fps u)$ the coefficient of $\mathcal Y^j$ in $\operatorname{pr}(X)$ lies in $u^{-1}\Bbbk\fps u$ for $j=1,\dots,4$ and in $\Bbbk\fps u$ for $j=0$.
\item[(g)] (Not used.) $x^5-\Bprod$ is separable, so $\tilde A$ is regular semisimple; in the situation of (f), $x^5-\Bprod$ is irreducible over $\Bbbk\lser u$, so $F[\mathcal Y]$ is a field.
\end{enumerate}
\end{lem}

\begin{proof}
(a) Induction on $m$, with $\Bprod^{(m)}_a\tilde b_{a+m}=\Bprod^{(m+1)}_a$ and $E_{a,a+5}=E_{aa}$; $\tilde A\sum_c\tilde b_c^{-1}E_{c+1,c}=1$ gives $\mathcal Y$, and $\mathcal Y^j$ follows likewise. Reindexing, $\mathcal Y^{5-m}=\sum_a(\Bprod^{(5-m)}_{a+m})^{-1}E_{a,a+m}$, and $\Bprod^{(m)}_a\Bprod^{(5-m)}_{a+m}=\Bprod$. In the Leibniz expansion of $\det(x1-\tilde A)$ only the identity and the even $5$-cycle $a\mapsto a+1$ contribute ($\pi(a)=a+1$ forces $\pi(a+1)=a+2$, and so on), the latter with $\prod_a(-\tilde b_a)=-\Bprod$.
(b) $\tilde AX=\sum_a\tilde b_a\xi_{a+1}E_{a,a+1+m}$ and $X\tilde A=\sum_a\xi_a\tilde b_{a+m}E_{a,a+m+1}$.
(c) By (b) $\operatorname{ad}\tilde A$ shifts the (direct) grading by $1$, so it suffices to treat $X\in\mathfrak{gl}^{(m)}$: $[\tilde A,X]=0$ means $\xi_{a+1}=\xi_a\tilde b_{a+m}/\tilde b_a$, whose solutions are determined by $\xi_0$, and $\xi_a=\Bprod^{(m)}_a$ is one. With $\tilde A^m=\Bprod\mathcal Y^{5-m}$ and $\mathcal Y=\Bprod^{-1}\tilde A^4$, $F[\tilde A]=F[\mathcal Y]$.
(d) For $Z\in\ker(\operatorname{ad}\tilde A)$, $\tr(Z[\tilde A,Y])=0$. For $X\in\mathfrak{gl}^{(m)}$, (a) gives $\tr(X\tilde A^{5-m})=\sum_a\xi_a\Bprod^{(5-m)}_{a+m}=\Bprod\lambda_m(X)$, so $\lambda_m$ vanishes on $\operatorname{ad}\tilde A(\mathfrak{gl}^{(m-1)})$, which has dimension $4$ by (c); as $\lambda_m(\tilde A^m)=5\neq0$, this image is $\ker\lambda_m$. Summing over $m$ gives the last assertion.
(e) $X=\frac15\lambda_m(X)\tilde A^m+\bigl(X-\frac15\lambda_m(X)\tilde A^m\bigr)$ with the second term in $\ker\lambda_m$; the second form uses (a). $X-\operatorname{pr}(X)\in\operatorname{im}(\operatorname{ad}\tilde A)$ is trace-orthogonal to $F[\mathcal Y]$ by (d), and the trace form is non-degenerate on $F[\tilde A]$: in the basis $1,\tilde A,\dots,\tilde A^4$ its Gram matrix has $5$ at $(0,0)$, $5\Bprod$ at $(i,5-i)$ and $0$ elsewhere ($\tilde A^n$ has shift $n$).
(f) The entries of $\mathcal Y$ are the $\tilde b_a^{-1}$, integral since $\tilde b_1^{-1}=u(u\tilde b_1)^{-1}$. The supports $\{(a+j,a)\}$ of $\mathcal Y^0,\dots,\mathcal Y^4$ partition the $25$ positions, and the entry $(a+j,a)$ of $\sum_jw_j\mathcal Y^j$ is $w_j(\Bprod^{(j)}_a)^{-1}$. For $j\leq4$ the window $\{a,\dots,a+j-1\}$ contains the pole index $1$ at most once, so $(\Bprod^{(j)}_a)^{-1}$ is integral, and a unit for $a=2$; this gives the criterion. $\Bbbk\fps u[\mathcal Y]$ is a ring since $\mathcal Y^5=\Bprod^{-1}1$, $\Bprod^{-1}\in u\Bbbk\fps u$, and by (c) it is the integral commutant. For integral $X\in\mathfrak{gl}^{(m)}$, $1\leq m\leq4$, the coefficient $\frac15\sum_a\xi_a\Bprod^{(5-m)}_{a+m}$ involves at most four consecutive weights, at most one of them $\tilde b_1$, so it lies in $u^{-1}\Bbbk\fps u$; for $m=5$ it is $\frac15\tr X$.
(g) $\gcd(x^5-\Bprod,5x^4)=1$. In the situation of (f), with $\Bprod^\circ:=u\Bprod$ (a unit) and $x=1/X$, $x^5-\Bprod=-\Bprod^\circ u^{-1}X^{-5}\bigl(X^5-u(\Bprod^\circ)^{-1}\bigr)$, an Eisenstein polynomial at $u$.
\end{proof}

\begin{lem}[Lifting]
\label{lem:gau-lift}
Let $\tilde A$ have the one-pole shape, let $X\in\mathfrak{gl}_5(\Bbbk\fps u)$ have shift $m+1$ with $0\leq m\leq4$, and suppose $\operatorname{pr}(X)=0$. Then there is $\mathcal Z\in\mathfrak{gl}_5(\Bbbk\fps u)$ of shift $m$ with $[\tilde A,\mathcal Z]=X$, and every solution $\mathcal Z'\in\mathfrak{gl}_5(\Bbbk\fps u)$ of $[\tilde A,\mathcal Z']=X$ is $\mathcal Z+w$ with $w\in\Bbbk\fps u[\mathcal Y]$. If $\tilde A$ and $X$ have rational coefficients, $\mathcal Z$ can be chosen rational.
\end{lem}

\begin{proof}
Write $X=\sum_a\xi_aE_{a,a+m+1}$ and try $\mathcal Z=\sum_a\rho_a\Bprod^{(m)}_aE_{a,a+m}$ (local letters). By \cref{lem:gau-comm}~(b), $[\tilde A,\mathcal Z]=\sum_a(\rho_{a+1}-\rho_a)\Bprod^{(m+1)}_aE_{a,a+m+1}$, so the equation is $\rho_{a+1}-\rho_a=\sigma_a:=\xi_a/\Bprod^{(m+1)}_a$, solvable if and only if $\sum_a\sigma_a=\lambda_{m+1}(X)=0$, i.e.\ $\operatorname{pr}(X)=0$ (\cref{lem:gau-comm}~(e)). Each $\sigma_a$ is integral, with $\sigma_a(0)=0$ when the window $\{a,\dots,a+m\}$ contains $1$. Put $\rho_1:=0$, $\rho_{a+1}:=\rho_a+\sigma_a$ ($a=1,\dots,4$); the relation for $a=0$ follows from $\sum_a\sigma_a=0$. The entry $\rho_a\Bprod^{(m)}_a$ is integral when the window $\{a,\dots,a+m-1\}$ avoids $1$, and otherwise ($a=2-m,\dots,1$) when $\rho_a(0)=0$: for $a=1$ by choice, and for $2-m\leq a\leq0$ because $\rho_a=-\sum_{c=a}^0\sigma_c$ with every window $\{c,\dots,c+m\}$ containing $1$. Two integral solutions differ by an element of $\Bbbk\fps u[\mathcal Y]$ (\cref{lem:gau-comm}~(c), (f)); the construction uses only field operations.
\end{proof}

\begin{lem}[The scalar equation]
\label{lem:gau-scal}
Let $A$, $\tilde A$ have the one-pole shape and satisfy (EP).
\begin{enumerate}
\item[(a)] For $w=\sum_{j=0}^4w_j\mathcal Y^j$ with $w_j\in\Bbbk\lser u$,
\[
\operatorname{pr}\bigl(w'+\operatorname{dlog}\Theta_U\,w\bigr)=\sum_{j=0}^4\bigl(w_j'+\chi^U_jw_j\bigr)\mathcal Y^j,\qquad \chi^U_j:=-\frac j5\,\frac{\Bprod'}{\Bprod}+\frac15\tr(\operatorname{dlog}\Theta_U).
\]
In particular $\operatorname{pr}((\mathcal Y^j)')=-\frac j5(\Bprod'/\Bprod)\mathcal Y^j$ and $\operatorname{pr}(\operatorname{dlog}\Theta_U\,\mathcal Y^j)=\frac15\tr(\operatorname{dlog}\Theta_U)\,\mathcal Y^j$. Moreover $\chi^U_j=\frac j{5u}+\chi^\circ_j$ with $\chi^\circ_j:=-\frac j5(\Bprod^\circ)'/\Bprod^\circ+\frac15\tr(\operatorname{dlog}\Theta_U)\in\Bbbk\fps u$, where $\Bprod^\circ:=u\Bprod\in\Bbbk\fps u^\times$.
\item[(b)] For $Z\in\mathfrak{gl}_5(\Bbbk\fps u)$ the equation $\operatorname{pr}(w'+\operatorname{dlog}\Theta_U\,w)=-\operatorname{pr}(Z)$ has solutions $w\in\Bbbk\fps u[\mathcal Y]$; their components $w_1,\dots,w_4$ are unique, and $w_0$ is unique up to $w_0\mapsto w_0+\cn\,\mathfrak d$, $\cn\in\Bbbk$.
\item[(c)] For $Z=0$ the solutions in $\Bbbk\fps u[\mathcal Y]$ are exactly the $\cn\,\mathfrak d\,1$, $\cn\in\Bbbk$.
\item[(d)] If the data are rational, so are the solutions for $\cn\in\bbQ$; $\mathfrak d\in1+u\bbQ\fps u$.
\end{enumerate}
\end{lem}

\begin{proof}
(a) $(\mathcal Y^j)'$ has shift $5-j$, with entries $((\Bprod^{(j)}_a)^{-1})'$ at $(a+j,a)$; by \cref{lem:gau-comm}~(e),
\[
\operatorname{pr}((\mathcal Y^j)')=\frac15\Bigl[\sum_a\bigl((\Bprod^{(j)}_a)^{-1}\bigr)'\Bprod^{(j)}_a\Bigr]\mathcal Y^j=-\frac15\Bigl[\sum_a\frac{(\Bprod^{(j)}_a)'}{\Bprod^{(j)}_a}\Bigr]\mathcal Y^j ,
\]
and $\sum_a(\Bprod^{(j)}_a)'/\Bprod^{(j)}_a=\sum_a\sum_{i<j}\tilde b_{a+i}'/\tilde b_{a+i}=j\Bprod'/\Bprod$, each index occurring $j$ times. Likewise $\operatorname{dlog}\Theta_U\mathcal Y^j$ has entries $\frac{(d^U_{a+j})'}{d^U_{a+j}}(\Bprod^{(j)}_a)^{-1}$, and \cref{lem:gau-comm}~(e) gives $\frac15\tr(\operatorname{dlog}\Theta_U)\mathcal Y^j$. As $\operatorname{pr}$ is $\Bbbk\lser u$-linear and fixes $\Bbbk\lser u[\mathcal Y]$, the formula follows. Finally $\Bprod'/\Bprod=(\Bprod^\circ)'/\Bprod^\circ-u^{-1}$, and $\tr(\operatorname{dlog}\Theta_U)$ is integral since the $d^U_c$ are units.

(b) By \cref{lem:gau-comm}~(f), $\operatorname{pr}(Z)=\sum_j\phi^Z_j\mathcal Y^j$ with $u\phi^Z_j\in\Bbbk\fps u$ for $j\geq1$ and $\phi^Z_0\in\Bbbk\fps u$ (local letters), and by (a) the equation splits into $w_j'+\chi^U_jw_j=-\phi^Z_j$. For $j=1,\dots,4$, multiplied by $u$, the coefficient of $u^n$ reads
\[
\Bigl(n+\frac j5\Bigr)w_{j,n}=-\coef{u^n}(u\phi^Z_j)-\sum_{l=0}^{n-1}\coef{u^{n-1-l}}\chi^\circ_j\cdot w_{j,l},
\]
with $n+j/5\neq0$: exactly one solution $w_j\in\Bbbk\fps u$. For $j=0$, $\chi^U_0=-\mathfrak d'/\mathfrak d$, so the equation is $\mathfrak d\,(w_0/\mathfrak d)'=-\phi^Z_0$, with solutions $w_0=\mathfrak d\bigl(\cn-\int_0^u\phi^Z_0/\mathfrak d\,\dd u\bigr)$, $\cn\in\Bbbk$ ($\int_0^u$ the formal antiderivative without constant term), and no others. (c) is (b) with $\phi^Z=0$. (d) The pivots are rational and $\mathfrak d$ is a binomial series of an element of $1+u\bbQ\fps u$.
\end{proof}

% Section 3, part 4 of 8: proof of Theorem U, the leading term, the Z/5-degree, graded gauges.

\begin{proof}[Proof of \cref{thm:gau-U}~(a)]
Write $P=\Theta_U\check P$ (a local letter); $\check P$ is integral if and only if $P$ is. As $A\Theta_U=\Theta_U\tilde A$, the gauge equation is equivalent to $z\check P'=[\tilde A,\check P]-z\operatorname{dlog}\Theta_U\,\check P$, i.e.\ to
\begin{equation}
\label{eq:gau-reduced}
[\tilde A,\check P_0]=0,\qquad [\tilde A,\check P_{k+1}]=\check P_k'+\operatorname{dlog}\Theta_U\,\check P_k\quad(k\geq0).
\end{equation}
We construct $\check P_k=\mathcal Z_k+w^{(k)}$ with $\mathcal Z_k$ integral and $w^{(k)}\in\Bbbk\fps u[\mathcal Y]$ such that (I$_k$) $[\tilde A,\check P_k]=\check P_{k-1}'+\operatorname{dlog}\Theta_U\check P_{k-1}$ ($\check P_{-1}:=0$) and (S$_k$) $\operatorname{pr}(\check P_k'+\operatorname{dlog}\Theta_U\check P_k)=0$. For $k=0$ take $\mathcal Z_0:=0$, $w^{(0)}:=\mathfrak d\,1$: (I$_0$) is clear and (S$_0$) is \cref{lem:gau-scal}~(c). Given $\check P_k$, $X_{k+1}:=\check P_k'+\operatorname{dlog}\Theta_U\check P_k$ is integral with $\operatorname{pr}(X_{k+1})=0$ by (S$_k$), so \cref{lem:gau-lift}, applied to each shift component, gives an integral $\mathcal Z_{k+1}$ with $[\tilde A,\mathcal Z_{k+1}]=X_{k+1}$; then \cref{lem:gau-scal}~(b) with $Z:=\mathcal Z_{k+1}'+\operatorname{dlog}\Theta_U\mathcal Z_{k+1}$ gives $w^{(k+1)}$, and (I$_{k+1}$) holds since $[\tilde A,w^{(k+1)}]=0$, (S$_{k+1}$) since its left side is $\operatorname{pr}(Z)+\operatorname{pr}\bigl((w^{(k+1)})'+\operatorname{dlog}\Theta_Uw^{(k+1)}\bigr)=0$. By \eqref{eq:gau-reduced}, $P^\natural:=\Theta_U\sum_kz^k\check P_k\in\Gau(A,\tilde A)$ and $P^\natural_0=\mathfrak d\,\Theta_U$.
\end{proof}

\begin{lem}[The leading term of a solution of the gauge equation]
\label{lem:gau-lead}
Let $A$, $\tilde A$ have the one-pole shape and satisfy (EP), let $P\in\Gau(A,\tilde A)$ and $n\geq0$ with $P_0=\dots=P_{n-1}=0$. Then $P_n=\cn\,\mathfrak d\,\Theta_U$ for a unique $\cn\in\Bbbk$, namely $\cn=(P_n)_{00}(0)$.
\end{lem}

\begin{proof}
For $\check P:=\Theta_U^{-1}P$, the coefficients of $z^n$ and $z^{n+1}$ of \eqref{eq:gau-reduced} give $[\tilde A,\check P_n]=0$, so $\check P_n\in\Bbbk\fps u[\mathcal Y]$ (\cref{lem:gau-comm}~(c), (f)), and $\check P_n'+\operatorname{dlog}\Theta_U\check P_n\in\operatorname{im}(\operatorname{ad}\tilde A)$, so its projection vanishes. By \cref{lem:gau-scal}~(c), $\check P_n=\cn\,\mathfrak d\,1$, and $\cn=(P_n)_{00}(0)$ as $d^U_0=\mathfrak d(0)=1$.
\end{proof}

\begin{proof}[Proof of \cref{thm:gau-U}~(b)--(d)]
Part (c) is \cref{prop:gau-graded}. (b) \emph{Injectivity:} if $\cn\neq0$ has order $n$, $\coef{z^n}(P^\natural\cn)=\cn_n\mathfrak d\Theta_U\neq0$. \emph{Surjectivity:} given $P\in\Gau(A,\tilde A)$, suppose $\cn_0,\dots,\cn_{n-1}$ are found with $R^{(n)}:=P-P^\natural\sum_{m<n}\cn_mz^m\equiv0$ modulo $z^n$. As $\Gau(A,\tilde A)$ is a $\Bbbk\fps z$-module, \cref{lem:gau-lead} gives $\coef{z^n}R^{(n)}=\cn_n\mathfrak d\Theta_U$ for a unique $\cn_n\in\Bbbk$, and then $R^{(n+1)}=R^{(n)}-\cn_nz^nP^\natural\equiv0$ modulo $z^{n+1}$. With $\cn:=\sum_n\cn_nz^n$, every coefficient of $P-P^\natural\cn$ vanishes. Its $z^0$-term is $\cn(0)\mathfrak d\Theta_U$; the statements on regular gauges follow, since $P^\natural\cn'=P\cdot(\cn_P^{-1}\cn')$ for $P=P^\natural\cn_P$ regular. For $(\tilde A,\tilde A)$, $\Theta_U=\mathfrak d=1$ and $1\in\Gau(\tilde A,\tilde A)$.

(d) Every object in the proof of (a) is obtained from the $b_a$, $\tilde b_a$ by field operations, derivatives and $\int_0^u$; with rational free constants the construction stays in $\mathfrak{gl}_5(\bbQ\fps u)\fps z$. Over $\Bbbk\supseteq\bbQ$ the proofs of \cref{lem:gau-comm,lem:gau-lift,lem:gau-scal,lem:gau-lead} hold verbatim, so (b) applies; and for rational $P=P^\natural\cn$ the recursion computes $\cn_n=(\coef{z^n}R^{(n)})_{00}(0)\in\bbQ$.
\end{proof}

\begin{lem}[The $\bbZ/5$-degree of the gauge equation]
\label{lem:gau-degree}
For $P=\sum_kz^kP_k$ with $P_k\in\mathfrak{gl}_5(\Bbbk\lser u)$ and $\delta\in\bbZ/5$ put $P^{[\delta]}:=\sum_kz^k\sum_{(a,b):\,k+b-a\equiv\delta}(P_k)_{ab}E_{ab}$, so that $P=\sum_\delta P^{[\delta]}$, and $P$ is graded if and only if $P=P^{[0]}$. Let $\operatorname{ge}(P):=zP'-AP+P\tilde A$. Then $\operatorname{ge}(P^{[\delta]})=\operatorname{ge}(P)^{[\delta+1]}$; in particular every degree part of a solution of the gauge equation is a solution.
\end{lem}

\begin{proof}
Give $z^kf(u)E_{ab}$ the degree $k+b-a$. Then $zP'$ raises it by $1$, and so do $A\,(fz^kE_{cd})=b_{c-1}fz^kE_{c-1,d}$ and $(fz^kE_{cd})\tilde A=\tilde b_dfz^kE_{c,d+1}$; so $\operatorname{ge}$ maps degree $\delta$ to degree $\delta+1$.
\end{proof}

\begin{prop}[Graded gauges]
\label{prop:gau-graded}
Let $A$, $\tilde A$ have the one-pole shape and satisfy (EP).
\begin{enumerate}
\item Graded regular gauges exist; they can be chosen rational if all $b_a$, $\tilde b_a$ are rational.
\item If $P$ is a graded regular gauge, the graded regular gauges are exactly the $P\cn(z)$ with $\cn\in\Bbbk\fps{z^5}^\times$.
\item For a graded regular gauge $P$, $P_{11}(0,z)=\sum_{m\geq0}(P_{5m})_{11}(0)\,z^{5m}$ is a unit of $\Bbbk\fps{z^5}$.
\item If $\zeta\in\Bbbk$, then $\operatorname{Rot}_iP=\sum_\delta\zeta^{-i\delta}P^{[\delta]}$; so $P$ is graded if and only if $\operatorname{Rot}_1P=P$, if and only if $\operatorname{Rot}_iP=P$ for all $i$, and $\frac15\sum_i\operatorname{Rot}_iP=P^{[0]}$. Moreover $A\,\zeta^{i(\mathrm{gr}+1)}=\zeta^i\zeta^{i(\mathrm{gr}+1)}A$, and the same for $\tilde A$; $\operatorname{Rot}_i$ maps solutions of $z\de_uY=AY$ to solutions, and $\operatorname{Rot}_i(PY)=(\operatorname{Rot}_iP)(\operatorname{Rot}_iY)$, $\operatorname{Rot}_i(Y\mathsf d(z))=\operatorname{Rot}_i(Y)\,\mathsf d(\zeta^{-i}z)$ for scalar series $\mathsf d$.
\end{enumerate}
\end{prop}

\begin{proof}
(i) By \cref{lem:gau-degree}, $(P^\natural)^{[0]}$ solves the gauge equation, and its $z^0$-term is the diagonal $P^\natural_0$; rationality is \cref{thm:gau-U}~(d). (ii) Regular gauges are $P\cn$, $\cn\in\Bbbk\fps z^\times$ (\cref{thm:gau-U}~(b)); with $\cn^{(r)}:=\sum_{n\equiv r}\cn_nz^n$, $(P\cn)^{[r]}=P\cn^{(r)}$, and $\cn\mapsto P\cn$ is injective, so $P\cn$ is graded if and only if $\cn\in\Bbbk\fps{z^5}$. (iii) For graded $P$, $(P_k)_{11}=0$ unless $5\mid k$, and $(P_0)_{11}(0)=\cn(0)\mathfrak d(0)d^U_1(0)\neq0$. (iv) The entry $(a,b)$ of $\coef{z^k}\operatorname{Rot}_iP$ is $\zeta^{-i(k+b-a)}(P_k)_{ab}$; the equivalences follow from $\sum_i\zeta^{-i\delta}=5\delta_{\delta,0}$. The entry $(a,a+1)$ of $\zeta^{i(\mathrm{gr}+1)}A\zeta^{-i(\mathrm{gr}+1)}$ is $\zeta^{-i}b_a$, so for a solution $Y$, $z\de_u(\operatorname{Rot}_iY)=\zeta^i\zeta^{i(\mathrm{gr}+1)}A\,Y(u,\zeta^{-i}z)=A\operatorname{Rot}_iY$. The product rules are immediate.
\end{proof}

% Section 3, part 5 of 8: (R1), the diagonal, the vanishing solutions.

%-------------------------------------------------------------------------------
\subsection{The residue relation, the vanishing solutions, (G) and symplecticity}
\label{sec:gau-van}

From now on, unless said otherwise, $A$ is the quintic system of \cref{prop:gau-quintic} in a conifold frame and $\tilde A$ is the model pulled back along $v=v(u)$; by \cref{prop:gau-equalprod} \cref{thm:gau-U} applies to them over $\bbQ$.

\begin{lem}[The residue relation (R1)]
\label{lem:gau-R1}
Let $A,\tilde A$ be cyclic systems with one simple pole over $\Bbbk\fps u$, with residues $r,\tilde r$, and let $P\in\Gau(A,\tilde A)$; put $P(0,z):=\sum_kz^kP_k(0)\in\mathfrak{gl}_5(\Bbbk\fps z)$. Then $rE_{12}P(0,z)=\tilde r\,P(0,z)E_{12}$: column $1$ of $P(0,z)$ is $P_{11}(0,z)\mathbf e_1$, row $2$ is $P_{22}(0,z)\mathbf e_2^T$, and $rP_{22}(0,z)=\tilde rP_{11}(0,z)$. If $r=\tilde r$, then $P_{22}(0,z)=P_{11}(0,z)$.
\end{lem}

\begin{proof}
$A=u^{-1}rE_{12}+O(1)$, likewise $\tilde A$, and $zP'$ has no $u^{-1}$; so the coefficient of $u^{-1}$ of the gauge equation is $rE_{12}P(0,z)=\tilde rP(0,z)E_{12}$, whose entry $(a,b)$ is $r\,\delta_{a1}P_{2b}(0,z)=\tilde r\,P_{a1}(0,z)\delta_{b2}$.
\end{proof}

\begin{lem}[The diagonal of the gauge]
\label{lem:gau-diag}
Let $A,\tilde A$ be cyclic systems with one simple pole over $\Bbbk\fps u$ with (EP), and $P$ a regular gauge. Then $P_0=\diag(p_{0,0},\dots,p_{0,4})$, and $\prod_ap_{0,a}=\cn(0)^5\det\Theta_U(0)=p_{0,1}(0)^5\prod_a\rho^U_a$ is a constant, where $\cn(0)\in\Bbbk^\times$ is as in \cref{thm:gau-U}~(b) and $\rho^U_a:=d^U_a(0)/d^U_1(0)$, i.e.\ $\rho^U_1=1$, $\rho^U_a=\prod_{c=1}^{a-1}(\tilde b_c/b_c)(0)$ for $a=2,3,4$ and $\rho^U_0=\prod_{c=1}^4(\tilde b_c/b_c)(0)$, with $(\tilde b_1/b_1)(0):=\tilde r/r$. Assume the normalisation
\[
\text{(N0)}\qquad b_a(0)=\tilde b_a(0)\ \ (a\neq1),\qquad r=\tilde r .
\]
Then $p_{0,a}(0)=p_{0,1}(0)$ for all $a$; if moreover $p_{0,1}(0)=1$, as for the canonical gauge of \cref{def:gau-can}, then $p_{0,a}(0)=1$ for all $a$, $\prod_ap_{0,a}=1$, and $p_{0,0}^5=\prod_{c=0}^3(b_c/\tilde b_c)^{4-c}$. For the quintic system in a conifold frame and the pulled-back model, where (N0) holds, and a regular gauge with $p_{0,1}(0)=1$ (for instance $\Pc$), this gives
\[
p_{0,0}(u)=\tau(u)^{-1/4}/\varpi_0,\qquad\text{i.e.}\qquad \tau^{1/4}p_{0,0}=1/\varpi_0 .
\]
\end{lem}

\begin{proof}
By \cref{thm:gau-U}~(b), $p_{0,a}=\cn(0)\mathfrak d\,d^U_a$, so $\prod_ap_{0,a}=\cn(0)^5\mathfrak d^5\det\Theta_U=\cn(0)^5\det\Theta_U(0)$ and $p_{0,a}(0)=p_{0,1}(0)\rho^U_a$; the recursion $d^U_{a+1}=d^U_a\tilde b_a/b_a$ at $u=0$ (round the cycle by (EP)) gives $\rho^U_a$. Under (N0) all $\rho^U_a=1$ and $d^U_a(0)=1$; if also $p_{0,1}(0)=1$, then $\cn(0)=1$, $\prod_ap_{0,a}=1$ and $p_{0,0}=\mathfrak d=(\det\Theta_U)^{-1/5}$ with $\det\Theta_U=\prod_{c=0}^3(\tilde b_c/b_c)^{4-c}$, as $c$ occurs in $4-c$ of the $d^U_a$. For the two systems of \cref{prop:gau-quintic,prop:gau-equalprod} the ratios $b_c/\tilde b_c$, $c=0,\dots,3$, are $\tc'/v'$, $v/(u\,v'\varpi_0^2(1-u)^3\tc'^2)$, $\tc'/v'$, $\varpi_0/((1-u)v')$, so $p_{0,0}^5=v'^{-10}v^3u^{-3}\varpi_0^{-5}(1-u)^{-10}$. By (EP), $v'^{-10}=u^2(1-u)^{10}v^{-2}$ (\cref{prop:gau-equalprod}~(ii)), hence $p_{0,0}^5=v/(u\varpi_0^5)=\tau^{-5/4}\varpi_0^{-5}$; both $p_{0,0}$ and $\tau^{-1/4}/\varpi_0$ lie in $1+u\bbQ\fps u$, where fifth roots are unique.
\end{proof}

\begin{defn}[Vanishing solutions]
\label{def:gau-van}
\begin{enumerate}
\item For $m\in\bbZ$ let $\IC_m(v,z):=\Res_{p=0}p^m\re^{-\Wv(p)/z}\dd p$ be the formal residue: in $\re^{-\Wv(p)/z}=\re^{-p^4/(4z)}\re^{-v/(pz)}$ the coefficient of each power of $v$ is a Laurent series in $p$ with coefficients in $\bbQ[z^{-1}]$, and $\Res_{p=0}$ takes its coefficient of $p^{-1}$. Explicitly
\[
\IC_m(v,z)=\sum_{i\geq0,\ m+4i+1\geq0}\frac{(-1/(4z))^i}{i!}\,\frac{(-v/z)^{m+4i+1}}{(m+4i+1)!}\in\bbQ[z^{-1}]\fps v .
\]
The \emph{vanishing solution of the model} is $\YC:=(\IC_0,-\IC_{-1},-\IC_3,\IC_2,\IC_1)$, i.e.\ $\YC_b=\varepsilon_b\IC_{\nu_b}$ (\cref{def:set-model}).
\item For a cyclic system $A=\sum_ab_aE_{a,a+1}$ with one simple pole over $\bbQ\fps u$ put $y^{(0)}:=-\mathbf e_1$, $y^{(n+1)}_a:=\int_0^ub_a\,y^{(n)}_{a+1}\,\dd u$ (the integrands are power series by \cref{lem:gau-vanq}) and $Y^{\mathrm C}_A:=\sum_{n\geq0}z^{-n}y^{(n)}$. For the quintic system, $\YCQ:=Y^{\mathrm C}_A$ is the \emph{vanishing solution of the quintic}.
\end{enumerate}
\end{defn}

\begin{lem}[The model's vanishing solution]
\label{lem:gau-vanmod}
(i) $z\,\de_v\IC_m=-\IC_{m-1}$; (ii) $mz\IC_{m-1}-\IC_{m+3}+v\IC_{m-2}=0$. Hence $z\,\dd\YC_b/\dd v=b^{\mathrm m}_b\YC_{b+1}$, $b^{\mathrm m}=(1,-v^{-1},1,-1,-1)$, and $\YC(0,z)=-\mathbf e_1$.
\end{lem}

\begin{proof}
(i) $\de_v\re^{-\Wv/z}=-(pz)^{-1}\re^{-\Wv/z}$. (ii) $\frac{\dd}{\dd p}\bigl(p^m\re^{-\Wv/z}\bigr)=\bigl(mp^{m-1}-z^{-1}(p^{m+3}-vp^{m-2})\bigr)\re^{-\Wv/z}$ has no residue. Then (i) gives $z\de_v\YC_b=b^{\mathrm m}_b\YC_{b+1}$ for $b\neq1$, and (ii) at $m=0$ gives $z\de_v\YC_1=\IC_{-2}=v^{-1}\IC_3=-v^{-1}\YC_2$. At $v=0$ only $m+4i+1=0$ survives, i.e.\ $m=-1$, $i=0$; so $\YC(0,z)=-\mathbf e_1$.
\end{proof}

\begin{lem}[Shape of $\YC$]
\label{lem:gau-vanshape}
For $n\geq0$ and $b\in\bbZ/5$, with $m:=\lfloor n/5\rfloor$,
\[
\coef{z^{-n}}\YC_b=\varepsilon_b\,(-1)^n\,\frac{v^{n-m}}{4^m\,m!\,(n-m)!}\quad\text{if }n\equiv\nu_b+1\pmod5,\qquad \coef{z^{-n}}\YC_b=0\ \text{otherwise}.
\]
Hence $\YC_b(v,z)=\varsigma^{\nu_b+1}y_b(z/\varsigma^4)$ with $y_b(t):=\YC_b(1,t)$ and $\varsigma=v^{1/5}$; the lowest terms are $\YC=(-v/z,\,-1,\,-v^4/(24z^4),\,-v^3/(6z^3),\,v^2/(2z^2))+\dots$; and at $v=v(u)$ the coefficient of $z^{-n}$ is $O(u^{n-m})$, so $\YC(v(u),z)\in\bbQ[z^{-1}]\fps u^5$ and $\YC(v(u),z)$ has no positive power of $z$.
\end{lem}

\begin{proof}
In \cref{def:gau-van}~(i) with $m'=\nu_b$ the term $i$ is a multiple of $v^{m'+4i+1}z^{-n}$, $n=m'+5i+1$; as $\nu_b+1\in\{0,\dots,4\}$, $i=m$, the sign is $(-1)^{5i+m'+1}=(-1)^n$, the exponent of $v$ is $n-m$ and the denominator $4^mm!\,(n-m)!$. Homogeneity follows from $v^{n-m}z^{-n}=\varsigma^{\nu_b+1}(z/\varsigma^4)^{-n}$, and $n-m\geq4n/5$ gives the last claim.
\end{proof}

\begin{lem}[Uniqueness of regular solutions]
\label{lem:gau-uniq}
Let $B=u^{-1}B_{-1}+\sum_{j\geq0}u^jB_j$ with $B_j\in\mathfrak{gl}_5(\Bbbk)$ independent of $z$. If $Y\in\Bbbk\lser z\fps u^5$ satisfies $z\de_uY=BY$ and $Y(0)=0$, then $Y=0$.
\end{lem}

\begin{proof}
For $Y=\sum_nu^nY^{(n)}$, the coefficient of $u^n$ is $\bigl((n+1)z-B_{-1}\bigr)Y^{(n+1)}=\sum_{j=0}^nB_jY^{(n-j)}$, and $\det((n+1)z\,1-B_{-1})$ is a polynomial in $z$ with leading term $(n+1)^5z^5$, a unit of $\Bbbk\lser z$; induct.
\end{proof}

\begin{lem}[The quintic's vanishing solution]
\label{lem:gau-vanq}
Let $A=\sum_ab_aE_{a,a+1}$ be a cyclic system with one simple pole over $\bbQ\fps u$, with residue $r$.
\begin{enumerate}
\item[(a)] $y^{(n)}_a=0$ unless $n\equiv1-a\pmod5$, and then $\ord_uy^{(n)}_a=(4n+\nu_a+1)/5$ exactly. If $b_a(0)=\tilde b_a(0)$ for $a\neq1$ and $r=\tilde r$, where $\tilde A$ is the pulled-back model, the leading coefficient of $y^{(n)}_a$ equals that of $\coef{z^{-n}}\YC_a(v(u),z)$.
\item[(b)] Every integrand $b_ay^{(n)}_{a+1}$ is a power series: the pole of $b_1$ meets $y^{(n)}_2$, which is non-zero only for $n\equiv4\pmod5$ and then has order $(4n+4)/5\geq4$. In particular no logarithm arises.
\item[(c)] $Y^{\mathrm C}_A\in\bbQ[z^{-1}]\fps u^5$, the coefficient of $z^{-n}$ having order at least $4n/5$; it solves $z\de_uY=AY$, and $Y^{\mathrm C}_A(0)=-\mathbf e_1$.
\item[(d)] $\coef{z^{-1}}Y^{\mathrm C}_A=-\bigl(\int_0^ub_0\,\dd u\bigr)\mathbf e_0$ and $\coef{u^0}(Y^{\mathrm C}_A)_0=0$. For the quintic system in a conifold frame, $\coef{z^{-1}}\YCQ=-\tc\,\mathbf e_0$.
\end{enumerate}
\end{lem}

\begin{proof}
Put $e(n,a):=(4n+\nu_a+1)/5$; for $n\equiv1-a$ it is a non-negative integer, as $\nu_a\equiv-a$ and $\nu_a+1\geq0$. Write $b_a=b^\circ_au^{-\delta_{a1}}(1+O(u))$ with $b^\circ_a:=b_a(0)$ for $a\neq1$ and $b^\circ_1:=r$ (local letters). We show by induction on $n$ that for $n\equiv1-a$, $y^{(n)}_a=c^{(n)}_au^{e(n,a)}(1+O(u))$ with $c^{(n)}_a\neq0$, and $y^{(n)}_a=0$ otherwise. For $n=0$, $y^{(0)}=-\mathbf e_1$ and $e(0,1)=0$. The component $y^{(n+1)}_a$ is built from $y^{(n)}_{a+1}$, which vanishes unless $n\equiv1-(a+1)$, i.e.\ $n+1\equiv1-a$. In that case, using $\nu_{a+1}=\nu_a-1+5\delta_{a1}$, the integrand is $b_ay^{(n)}_{a+1}=b^\circ_ac^{(n)}_{a+1}u^{e(n+1,a)-1}(1+O(u))$ with $e(n+1,a)-1=(4n+\nu_a)/5\geq0$, a power series, and integration gives $c^{(n+1)}_a=b^\circ_ac^{(n)}_{a+1}/e(n+1,a)\neq0$. This proves the first part of (a) and (b); for $a=1$ the integrand involves $y^{(n)}_2$ with $n\equiv4$, of order $e(n,2)=(4n+4)/5\geq4$. The pulled-back model satisfies the same recursion for the leading coefficients with $\tilde b^\circ_a$ in place of $b^\circ_a$, so under the stated hypothesis the leading coefficients of $Y^{\mathrm C}_A$ and $Y^{\mathrm C}_{\tilde A}$ agree; and $Y^{\mathrm C}_{\tilde A}=\YC(v(u),z)$ by \cref{lem:gau-uniq}, both being solutions of $z\de_uY=\tilde AY$ with value $-\mathbf e_1$ (\cref{lem:gau-vanmod} and (c) below for $\tilde A$). (c) The order bound is $e(n,a)\geq4n/5$, so each power of $u$ receives finitely many $n$. The definition says $(y^{(n+1)})'=Ay^{(n)}$, hence $z\de_uY^{\mathrm C}_A=\sum_{n\geq1}z^{1-n}Ay^{(n-1)}=AY^{\mathrm C}_A$; and $y^{(n)}(0)=0$ for $n\geq1$. (d) $y^{(1)}_a=-\delta_{a0}\int_0^ub_0$, and $(Y^{\mathrm C}_A)_0(0)=y^{(0)}_0=0$. For the quintic, $b_0=\tc'$ and $\tc(0)=0$. In a conifold frame the hypothesis of (a) holds (\cref{prop:gau-equalprod}).
\end{proof}

% Section 3, part 6 of 8: the gauge identity (G), the canonical gauge, symplecticity, HPF.

\begin{thm}[The gauge identity (G)]
\label{thm:gau-G}
For every $P\in\Gau(A,\tilde A)$, in particular for every regular gauge from the pulled-back model to the quintic system,
\begin{equation}
\label{eq:gau-G}
P(u,z)\,\YC(v(u),z)=P_{11}(0,z)\,\YCQ(u,z)\qquad\text{in }\bbQ\lser z\fps u^5 .
\end{equation}
\begin{enumerate}
\item[(a)] For a graded regular gauge $P$, $P_{11}(0,z)$ is a unit of $\bbQ\fps{z^5}$, $\Pc:=P\,P_{11}(0,z)^{-1}$ is a graded regular gauge, and $\Pc$ is the unique regular gauge $P'$ with $P'\,\YC(v(u),z)=\YCQ$.
\item[(b)] For a regular gauge $P$, $P\,\YC(v(u),z)$ has no positive power of $z$ if and only if $P_{11}(0,z)$ is a constant; the regular gauges with this property are the $\mu\Pc$, $\mu\in\bbQ^\times$.
\item[(c)] For any two cyclic systems $A,\tilde A$ with one simple pole over $\bbQ\fps u$ and every $P\in\Gau(A,\tilde A)$: $P\,Y^{\mathrm C}_{\tilde A}=P_{11}(0,z)\,Y^{\mathrm C}_A$.
\end{enumerate}
\end{thm}

\begin{proof}
We prove (c), of which \eqref{eq:gau-G} is the case of the two systems of this subsection: $\YC(v(u),z)=Y^{\mathrm C}_{\tilde A}$ by \cref{lem:gau-uniq}, both being solutions of $z\de_uY=\tilde AY$ with value $-\mathbf e_1$ (\cref{lem:gau-vanmod,lem:gau-vanq}~(c)). Put $\tilde Y:=Y^{\mathrm C}_{\tilde A}\in\bbQ[z^{-1}]\fps u^5$. The product $P\tilde Y$ lies in $\bbQ\lser z\fps u^5$: at each order in $u$ it is a finite sum of products of a power series in $z$ and a polynomial in $z^{-1}$. It solves the first system: $z\de_u(P\tilde Y)=(AP-P\tilde A)\tilde Y+P\tilde A\tilde Y=A(P\tilde Y)$. Its value at $u=0$ is $P(0,z)(-\mathbf e_1)=-P_{11}(0,z)\mathbf e_1$ by \cref{lem:gau-R1}. $P_{11}(0,z)Y^{\mathrm C}_A$ is also a solution, $P_{11}(0,z)$ being constant in $u$, with the same value by \cref{lem:gau-vanq}~(c). \Cref{lem:gau-uniq} with $B=A$ gives the identity.

(a) $P_{11}(0,z)$ is a unit of $\bbQ\fps{z^5}$ by \cref{prop:gau-graded}~(iii), so $\Pc$ is a graded regular gauge by \cref{prop:gau-graded}~(ii), and $(\Pc)_{11}(0,z)=1$; by \eqref{eq:gau-G}, $\Pc\YC(v(u))=\YCQ$. Conversely, let $P'$ be any regular gauge with $P'\YC(v(u))=\YCQ$. By \cref{thm:gau-U}~(b), (d), $P'=\Pc\cn(z)$ for some $\cn\in\bbQ\fps z^\times$, so that $\YCQ=\cn(z)\Pc\YC(v(u))=\cn(z)\YCQ$, i.e.\ $(\cn(z)-1)\YCQ=0$ in $\bbQ\lser z\fps u^5$; the coefficient of $u^0$ is $-(\cn(z)-1)\mathbf e_1$, so $\cn=1$. (b) If $P_{11}(0,z)=\mu$ is constant, then $P\YC(v(u))=\mu\YCQ$ lies in $\bbQ[z^{-1}]\fps u^5$. Conversely, the coefficient of $u^0$ of $P\YC(v(u))$ is $-P_{11}(0,z)\mathbf e_1$ with $P_{11}(0,z)\in\bbQ\fps z$, so it has no positive power of $z$ only if $P_{11}(0,z)$ is constant. With $P=\Pc\cn(z)$ we have $P_{11}(0,z)=\cn(z)$, which gives the last assertion.
\end{proof}

\begin{defn}[The canonical gauge]
\label{def:gau-can}
In a conifold frame, the \emph{canonical gauge} $\Pc$ is the regular gauge from the pulled-back model to the quintic system with $(\Pc)_{11}(0,z)=1$. It exists and is unique: $\Pc=P\,P_{11}(0,z)^{-1}$ for any graded regular gauge $P$ (\cref{thm:gau-G}~(a)), and a regular gauge $\Pc\cn(z)$ (\cref{thm:gau-U}~(b)) has $(\Pc\cn)_{11}(0,z)=\cn(z)$. Equivalently, $\Pc$ is the unique regular gauge with $\Pc\YC(v(u))=\YCQ$. It is graded, rational (\cref{thm:gau-U}~(d)), colour covariant, $\operatorname{Rot}_i\Pc=\Pc$ (\cref{prop:gau-graded}~(iv)), and $p_{0,a}(0)=1$ for its diagonal $z^0$-term (\cref{lem:gau-diag}). We write $\Pc=\sum_kz^k\sum_bp_{k,b}(u)E_{b+k,b}$.
\end{defn}

\begin{prop}[Symplecticity]
\label{prop:gau-symp}
Let $A,\tilde A$ be self-dual cyclic systems with one simple pole over $\Bbbk\fps u$ with (EP) and $r=\tilde r$, and let $P$ be a regular gauge from $\tilde A$ to $A$, $\cn(z):=P_{11}(0,z)$. Then
\begin{equation}
\label{eq:gau-symp}
P(u,-z)^T\eta\,P(u,z)=\cn(z)\cn(-z)\,\eta .
\end{equation}
In particular, for the two systems of \cref{sec:gau-van}, $\Pc$ is symplectic, $\Pc(u,-z)^T\eta\Pc(u,z)=\eta$; the symplectic regular gauges are the $\Pc\cn(z)$ with $\cn(z)\cn(-z)=1$, so symplecticity constrains only the even part of $\log(\cn/\cn(0))$; for graded gauges, $\cn=\pm\exp(a_1z^5+a_3z^{15}+a_5z^{25}+\cdots)$. Among the symplectic regular gauges, only (G) singles out $\Pc$.
\end{prop}

\begin{proof}
For $S:=P(u,-z)^T\eta P(u,z)$, the gauge equation at $-z$, transposed, and self-duality $A^T\eta=\eta A$, $\tilde A^T\eta=\eta\tilde A$ (\cref{def:gau-onepole}) give
\[
zS'=\bigl(-P(-z)^TA^T+\tilde A^TP(-z)^T\bigr)\eta P(z)+P(-z)^T\eta\bigl(AP(z)-P(z)\tilde A\bigr)=\tilde A^TS-S\tilde A,
\]
so $T:=\eta^{-1}S$ satisfies $zT'=\tilde AT-T\tilde A$ and $T=t(z)\,1$ by \cref{thm:gau-U}~(b) for $(\tilde A,\tilde A)$. At $u=0$, by \cref{lem:gau-R1} (column $1$ of $P(0,z)$ is $\cn(z)\mathbf e_1$, and $P_{22}(0,\cdot)=\cn$ as $r=\tilde r$) and $\eta_{a1}=5\delta_{a2}$, $S_{21}(0,z)=5P_{22}(0,-z)\cn(z)=5\cn(-z)\cn(z)$, while $S_{21}=5t(z)$. The consequences follow from $(\Pc\cn)_{11}(0,z)=\cn(z)$: $\cn(z)\cn(-z)=1$ forces $\cn(0)=\pm1$ and is equivalent to the oddness of $\log(\cn/\cn(0))$, and graded gauges have $\cn\in\Bbbk\fps{z^5}$.
\end{proof}

%-------------------------------------------------------------------------------
\subsection{The normalisation series}
\label{sec:gau-R0}

We compare $\Pc$ applied to the model columns with the rows of \cref{def:set-rows}: colour $0$ in \cref{thm:gau-R0}, the others by colour covariance (\cref{cor:gau-colours}). Here $\theta_u:=u\,\dd/\dd u$, $\theta_{\zh}:=\zh\,\dd/\dd\zh$, and a series $\Phi(\zh,u)=\sum_n\zh^n\mathrm{ph}_n(u)$ with $\mathrm{ph}_n\in\bbQ\lser s$ is identified with $\sum_nz^ns^{-4n}\mathrm{ph}_n$ through $\zh=z/s^4=L^4z$ (\cref{def:set-points}); $\Dq$ acts on $\bbQ\lser s$ as in \cref{lem:gau-confexp}.

\begin{lem}[HPF]
\label{lem:gau-HPF}
For $e\in\bbZ$ let $\mathscr O_e:=1+\zh(1-u)\bigl(\tfrac45\theta_{\zh}-\tfrac e5-\theta_u\bigr)$, and for $\Phi=\sum_n\zh^n\mathrm{ph}_n$ as above put $\Phi^{[0]}:=\Phi$ and $\Phi^{[j]}:=\mathscr O_{-j}\Phi^{[j-1]}$, $j=1,\dots,5$.
\begin{enumerate}
\item $L\Phi$ satisfies $\PFfive^{(0)}(L\Phi)=0$, where $\PFfive^{(0)}$ is $\PFfive$ with $z\Dq$ replaced by $z\Dq+L$ (that is, $\re^{u_0/z}L\Phi$ solves $\PFfive$ formally), if and only if
\begin{equation}
\label{eq:gau-HPF}
\Phi^{[5]}-\Phi-(1-u)\Bigl[3\zh\Phi^{[4]}+\tfrac{17}5u\zh^2\Phi^{[3]}+\tfrac95u^2\zh^3\Phi^{[2]}+\tfrac{274}{625}u^3\zh^4\Phi^{[1]}+\tfrac{24}{625}u^4\zh^5\Phi\Bigr]=0 .
\end{equation}
\item The coefficient of $\zh^{n+1}$ of the left side of \eqref{eq:gau-HPF} is $(1-u)\bigl[(4n-5\theta_u)\mathrm{ph}_n+r_n\bigr]$, where, with $\Phi^\flat_{j,n}$ the coefficient of $\zh^n$ of $\Phi^{[j]}$ computed with $\mathrm{ph}_n$ replaced by $0$ (and $\Phi^{[j]}_m:=\coef{\zh^m}\Phi^{[j]}$, zero for $m<0$),
\[
r_n:=\sum_{j=1}^5\Bigl(\frac{4n+j}5-\theta_u\Bigr)\Phi^\flat_{j-1,n}-3\Phi^\flat_{4,n}-\tfrac{17}5u\Phi^{[3]}_{n-1}-\tfrac95u^2\Phi^{[2]}_{n-2}-\tfrac{274}{625}u^3\Phi^{[1]}_{n-3}-\tfrac{24}{625}u^4\mathrm{ph}_{n-4}
\]
depends only on $\mathrm{ph}_0,\dots,\mathrm{ph}_{n-1}$. So \eqref{eq:gau-HPF} holds if and only if $(4n-5\theta_u)\mathrm{ph}_n=-r_n$ for all $n\geq0$; here $r_0=0$, and $r_1=\frac35$ when $\mathrm{ph}_0=1$.
\item If every $\mathrm{ph}_n$ lies in $\bbQ\fps u$, \eqref{eq:gau-HPF} reads $(4n-5k)\coef{u^k}\mathrm{ph}_n=-\coef{u^k}r_n$ for all $n,k\geq0$: for $5k\neq4n$ the coefficient $\coef{u^k}\mathrm{ph}_n$ is fixed by $\mathrm{ph}_0,\dots,\mathrm{ph}_{n-1}$, while for $n=5m$ the coefficient $\fc_m:=\coef{u^{4m}}\mathrm{ph}_{5m}$ is free and $\coef{u^{4m}}r_{5m}=0$ is required. In particular $\mathrm{ph}_0$ is a constant, and $\mathrm{ph}_0=1$ forces $\mathrm{ph}_1=-\frac3{20}$.
\end{enumerate}
\end{lem}

\begin{proof}
(1) From $\Dq(s^cf)=(1-u)s^{c-5}\bigl(-\frac c5-\theta_u\bigr)f$ (as $\Dq=-(1-u)u^{-1}\theta_u$), the coefficient of $z^m$ of $(z\Dq+L)(s^eF^\flat)$, $F^\flat=\sum_m\zh^mF^\flat_m$ (a local letter), is $s^{e-4m-1}\bigl[F^\flat_m+(1-u)\bigl(\frac{4(m-1)-e}5-\theta_u\bigr)F^\flat_{m-1}\bigr]$; so $(z\Dq+L)(s^eF^\flat)=s^{e-1}\mathscr O_eF^\flat$. (2) By induction $(z\Dq+L)^j(s^{-1}\Phi)=s^{-1-j}\Phi^{[j]}$; as $z=s^4\zh$, $z^m(z\Dq+L)^{5-m}(L\Phi)$ is $s^{-1}u^{m-1}\zh^m\Phi^{[5-m]}$. (3) With $M=z\Dq$, $5^5q=1-u$ and the elementary symmetric functions $(\sigma_1,\dots,\sigma_5)=(3,\frac{17}5,\frac95,\frac{274}{625},\frac{24}{625})$ of $\frac15,\dots,\frac55$, $\PFfive=uM^5-1-(1-u)\sum_{k=1}^5\sigma_kz^kM^{5-k}$; replacing $M$ by $z\Dq+L$, applying to $L\Phi$ and using (2) gives $s^{-1}$ times the left side of \eqref{eq:gau-HPF}, which is (i).
(4) By definition $\Phi^{[j]}_m=\Phi^{[j-1]}_m+(1-u)\bigl(\frac{4(m-1)+j}5-\theta_u\bigr)\Phi^{[j-1]}_{m-1}$, so $\Phi^{[5]}_{n+1}-\mathrm{ph}_{n+1}=(1-u)\sum_{j=1}^5(\frac{4n+j}5-\theta_u)\Phi^{[j-1]}_n$; with $\Phi^{[j]}_n=\mathrm{ph}_n+\Phi^\flat_{j,n}$ the terms in $\mathrm{ph}_n$ are $(1-u)\bigl[\sum_{j=1}^5(\frac{4n+j}5-\theta_u)-3\bigr]\mathrm{ph}_n=(1-u)(4n-5\theta_u)\mathrm{ph}_n$, and the rest is $(1-u)r_n$. Clearly $r_0=0$; for $\mathrm{ph}_0=1$, $\Phi^\flat_{j,1}=(1-u)\frac{j(j+1)}{10}$ and $(c-\theta_u)(1-u)=c(1-u)+u$ give $r_1=\sum_{j=2}^5\frac{j(j-1)}{10}\bigl(\frac{4+j}5(1-u)+u\bigr)-6(1-u)-\frac{17}5u=\frac35$.
(5) (iii) follows from $\theta_uu^k=ku^k$: for $n=0$ it gives $\mathrm{ph}_0\in\bbQ$, and for $n=1$, with $r_1=\frac35$, $\mathrm{ph}_1=-\frac3{20}$.
\end{proof}

\input{sections/s03_gauge_7}



